\documentclass[10pt,a4paper]{amsart}
\usepackage[margin=19mm]{geometry}
\usepackage{amsmath,amssymb,mathtools}
\usepackage[scr=rsfs,cal=euler]{mathalfa}
\usepackage{xfrac}
\usepackage[unicode,hidelinks,bookmarks=false]{hyperref}
\newcommand{\ie}{i.e.\ }
\newcommand{\cf}{cf.\ }
\newcommand{\resp}{respectively\ }
\newcommand{\Subsection}{\subsection}
\providecommand{\nobreakdash}{\nobreak\hbox{-}}
\setlength{\emergencystretch}{2em}
\multlinegap0pt
\usepackage{amssymb}
\usepackage{mathtools}
\newtheorem{lemma}{Lemma}
\newtheorem{theorem}{Theorem}
\newcommand{\HOX}[1]{\todo[noline,color=white,size=\footnotesize]{#1}}
\def\R{\mathbb R}
\title{Introduction to inverse problems for hyperbolic PDEs}
\author{Medet Nursultanov, Lauri Oksanen}
\date{Source: arXiv:2311.06020v2 (CC BY 4.0)}
\begin{document}
\maketitle

\noindent\emph{Source and licence.} Introduction to inverse problems for hyperbolic PDEs. Version arXiv:2311.06020v2 (CC BY 4.0), distributed by arXiv under the Creative Commons Attribution 4.0 International licence, http://creativecommons.org/licenses/by/4.0/, as recorded in the arXiv licence metadata for this exact version and shown on the abstract page https://arxiv.org/abs/2311.06020v2. Full licence notice: \texttt{licenses/arxiv-2311.06020-CC-BY-4.0.txt}.\par\smallskip

\noindent\textit{Editorial scope.} Counted unit: the article's theorem th\_1d\_main (source0.tex lines 537--539). The article's complete original proof of it is delivered with it (source0.tex lines 540--593, one \texttt{proof} environment of 54 lines). No mathematics is written, repaired or re-derived: the statement and the proof are the article's own lines, in the article's own notation, and no retained line is substituted or re-flowed.  Retained uncounted as the source-local context the proof needs: lines 511--514; lines 531--533.  Nothing was omitted from any retained span, so the package carries no omission record.  No retained line contains a citation, so the wrapper carries no bibliography and no bibliography entry.  No retained line contains a figure environment, an image-inclusion command or drawing code, so no image is delivered and the package declares no asset dependency.  Editorial formatting only, in addition: an AMS \texttt{amsart} preamble that restates the article's own declarations and macros used by the retained lines, namely the environments \texttt{lemma}, \texttt{theorem} on separate counters, together with the standard packages those lines need.  The retained environments print in this document as Lemma 1, Theorem 1 in order of appearance, whereas the article prints the same items with the numbering of its own full document.  The complete native source of the article as distributed in the arXiv e-print for this exact version is bundled unchanged as \texttt{sources/149825/source0.tex}.
	\begin{align}\label{scalar prod supp 0}
		\left(1_{(0,s)}u_1^f (T, \cdot), u_1^h (T, \cdot)\right)_{L^2(0,1)} =
		 \left(1_{(0,s)}u_2^f (T, \cdot), u_2^h (T, \cdot)\right)_{L^2(0,1)}.
	\end{align}

\begin{lemma}\label{lem_nonvanishing}
	Let $T\geq 1$, then for any $x_0\in (0,1)$ there is $f\in C_0^\infty(0,T)$ such that $u^f(T,x_0) \neq 0$. 
\end{lemma}

\begin{theorem}\label{th_1d_main}
	If $\Lambda_1=\Lambda_2$, then $q_1=q_2$.
\end{theorem}

\begin{proof}
We consider $T > 1$ and let $x\in (0,1)$. Let us differentiate the left-hand side of \eqref{scalar prod supp 0}:
%\HOX{At this point we rename $s$ to $x$, see the previous comment}
\begin{multline*}
	\partial_x 	\left(1_{(0,x)}(\cdot)u_1^f (T, \cdot), u_1^h (T, \cdot)\right)_{L^2(0,1)} \\
	= \partial_x\int_{0}^{x} u_1^f (T, y) u_1^h (T, y) dy = u_1^f (T, x) u_1^h (T, x).
\end{multline*}
Therefore, by \eqref{scalar prod supp 0}, we obtain
\begin{equation}\label{bububu}
	u_1^f (T, x) u_1^h (T, x) = u_2^f (T, x) u_2^h (T, x).
\end{equation}

Due to Lemma \ref{lem_nonvanishing},
for each $x \in (0,1)$ there is $f \in C_0^\infty(\R^+)$ such that 
$u_1^f(T,x) \ne 0$. Choosing such $f$ we may define
	\begin{align*}
w(x) = \frac{u_2^f (T, x)}{u_1^f (T, x)}.
	\end{align*}
We emphasize that the choice of $f$ depends on $x$, and it may appear that $w$ could be non-smooth. However, it is smooth. Indeed, 
	\begin{align}\label{w_eq}
u_1^h (T, x) = w(x) u_2^h (T, x),
	\end{align}
and for each $x_0 \in (0,1)$ there is $h \in C_0^\infty(\R^+)$ such that 
$u_2^h(T,x_0) \ne 0$. Thus, for $x$ near $x_0$,
	\begin{align*}
w(x) = \frac{u_1^h (T, x)}{u_2^h (T, x)}.
	\end{align*}
As the right-hand side is smooth for $x$ near $x_0$, and as $x_0 \in (0,1)$ is arbitrary, we see that $w$ is smooth. 

Let us now take $f = h$ in \eqref{bububu} and use \eqref{w_eq},
	\begin{align*}
(u_2^h(T,x))^2 = (u_1^h(T,x))^2 = w^2(x) (u_2^h(T,x))^2.
	\end{align*} 
Choosing again $x \in (0,1)$ and $h \in C_0^\infty(\R^+)$ such that 
$u_2^h(T,x) \ne 0$, we see that $w^2(x) = 1$.
The smoothness of $w$ implies that it is a constant function
taking the value $1$ or $-1$. 

To summarize 
	\begin{align*}
u_1^h (T, x) = \pm u_2^h (T, x),
	\end{align*}
for all $T > 1$, $x \in (0,1)$ and $h \in C_0^\infty(\R^+)$.


There holds
	\begin{align*}
0 = (\partial_t^2 - \partial_x^2 + q_1) u_1^h 
= \pm (\partial_t^2 - \partial_x^2 + q_1) u_2^h 
= \pm (q_1 - q_2) u_2^h.
	\end{align*}
Choosing such $h \in C^\infty_0(\R^+)$ that $u_2^h(T,x) \ne 0$, we get 
$q_1(x) = q_2(x)$.
\end{proof}

\end{document}
